\section{Átomos del conteo y diferencias en las supresiones}
Para el logaritmo exacto aproximado se suman los conjuntos independientes $X$ de sitios
sobre el exterior común. El grafo de incompatibilidad une sitios en las fibras RC y
diagonales activadas incompatibles. Si la activación sólo prohíbe ciertas parejas
dentro de una fibra, el mismo argumento usa ese subgrafo. Las actividades son las de
la silla local. El peso sobre $X$, aparte de $\zeta^X$, es la exponencial de
$G_s(X)=R_c(X)+A_c(X)+\widetilde d_T(X)$, normalizada en $X=\varnothing$.
$A_c$ es el prefactor exterior centrado y $\widetilde d_T$ la diferencia de errores
de torres con aproximación común de orden $q=4v+12$.
\paragraph{Dependencia del conjunto suprimido.}
En un átomo RC compatible cada fila y columna escogidas se suprimen una sola vez.
Ponemos $T_X=\sum_{i\in X}(1,\delta_{r_i},\delta_{c_i})$ y $\theta=sT_X$.
Las dos medidas supervivientes son $\mu^0-\theta_R,\mu^0-\theta_C$, de igual masa
$1-\theta_0$. Hasta constante y un campo afín de coeficientes acotados, $G_s(X)$ es
\[
       s^{-1}\mathscr A(s,\theta)+\mathscr B_q(s,\theta).
\]
Ambas funciones son holomorfas acotadas en una bola de TV fija en el subespacio de
masas compatibles, con $|s|\le\kappa/v$. Los campos afines inicialmente no acotados
ya se extrajeron en $\zeta$. El centrado restante de $R_c$ tiene gradiente acotado
por el apartado anterior. Detallamos por qué la bola de medidas tiene radio fijo.
\begin{itemize}
\item En la expresión integral de $\mathcal R$, escrita en $w,M_1$ tras cancelar
$t\log t$, sólo aparecen los supervivientes, con $t=1$. Se resuelve el equilibrio con
esas medidas en la bola por la misma brecha. En $x=\lambda\mathbf1_X$,
$0\le\lambda\le1$, los $t$ valen $1$ o $1-\lambda$ y la perturbación de medidas es
$O(v|s|)$ hasta el extremo. La expresión ya cancelada admite el valor en el átomo.
\item En el prefactor, con $\sigma$ independiente de la medida,
$l-w=\sigma t_rt_c w^2/(1-2\sigma t_rt_c w)$. El margen da grado $1+O(\sigma)$
uniformemente y grado uno en los índices apagados. El logaritmo del grado dividido
por $\sigma$ es acotado en un disco fijo. También lo es
$\log b_2(\sigma t_rt_c w)/\sigma^2$ por $\log b_2(y)=-y^2+O(y^3)$.
El cofactor centrado usa transposición bilineal. La serie de Fredholm está dominada
por $\sum_l C^l l^{l/2}/l!$, y la brecha basal separa de cero el determinante.
Se conservan los términos de masa del cofactor al pasar de $j$ a $m=j-|X|$,
como en \emph{Media matriz}, parte 05. Dependen analíticamente de $1-\theta_0$.
Todo esto vale en el polidisco local antes de ir al átomo con medidas $t\mu$.
\item Para el error puro de torres se normalizan las medidas por $1-\theta_0$ y
se balancea el mismo núcleo. La parte 03 de \emph{Grados intermedios} da la función
$G_q^{\rm tor}(s_m,z_X)$ acotada en una bola uniforme de datos y $|s_m|\le c/q$.
Aquí $s_m=s/(1-\theta_0)$, y basta reducir $\kappa$. En el menor real la cota es
$\Delta_q=(Cq/j)^q+Ce^{-cj}$, ajustando las constantes. Se toma la diferencia al
padre. El factor inicial $s_m$ y las anulaciones por grupos de esa entrada se conservan.
\end{itemize}
En conjuntos incompatibles extendemos las funciones por la misma fórmula en medidas
con signo. Esta extensión sirve sólo para diferencias de funciones. Los términos
físicos siguen restringidos a los conjuntos compatibles.
\paragraph{Diferencias mixtas.}
Sea $\Delta_iG(Y)=G(Y\cup\{i\})-G(Y)$ para $i\notin Y$ y sea $\Delta_D$ el
producto de diferencias en índices distintos. Cada dirección $T_i$ tiene norma
a lo sumo tres en la suma de las tres normas. Con margen fijo a la frontera de la
bola quedan dentro todas las interpolaciones de supresiones, pues $3v|s|\le3\kappa$.
Cauchy para las derivadas multilineales de orden $l+1$ da $(l+1)!C^{l+1}$:
se pueden tomar radios iguales a una constante dividida por $l+1$ y usar Stirling.
Integrando esas derivadas en las direcciones $sT_i$, se obtiene
\[
 |\Delta_iG_s(Y)|\le B_0,\qquad
 |\Delta_i\Delta_DG_s(Y)|\le(l+1)! C^{l+1}|s|^l,\quad |D|=l\ge1 .
\]
El polo de $\mathscr A/s$ se cancela al normalizar en el vacío. En particular,
para cualquier conjunto restante $W$ disjunto de $Y,i$,
$\sum_{\varnothing\ne D\subseteq W}|\Delta_D\Delta_i G_s(Y)|
\le C\sum_{l\ge1}(l+1)(Cv|s|)^l\le B_1$, uniformemente en $v$.
